% Dodatak A: Šabloni promptova i primeri ocenjivanja
\chapter{Šabloni promptova i primeri ocenjivanja}
\label{app:promptovi}

Dodatak sadrži šablone promptova korišćene u eksperimentima i nekoliko primera ocenjivanja.
Promptovi su preneti doslovno iz konfiguracionih datoteka, sa izvornim nedoslednostima u dijakritici, jer su modeli primili takav tekst.

\medskip
\noindent\textbf{Napomena.} Model ocenjuje na skali 0--100,
a ocene se u analizi i u primerima ispod preslikavaju na skalu 0--10, radi poređenja sa ocenama nastavnika.

\section*{Uputstvo o formatu izlaza}

Sledeće uputstvo dodaje se na kraj sistemske instrukcije u svim konfiguracijama strogosti i u svim režimima ocenjivanja:

\begin{tcolorbox}[promptbox]

Tvoj odgovor treba da bude ocena studentskog odgovora u sledećem formatu:

Prvi red treba da sadrži samo broj na skali od 0 do 100, gde je 100 najbolja ocena a 0 najgora.

Ostatak odgovora treba da sadrži obrazloženje ocene.

\end{tcolorbox}

\section*{Promptovi strogosti za PRA režim}

U PRA režimu prompt uz sistemsku instrukciju sadrži pitanje, referentni odgovor nastavnika (označen kao „tačan odgovor'')
i studentski odgovor, u zasebnim blokovima.

\begin{tcolorbox}[promptbox, title={Blaga konfiguracija}]

Ti si blag profesor visokog obrazovanja i treba da ocenis odgovore studenata.

Dobijas pitanje, tacan odgovor kao i odgovor studenta, i treba da proveris da li je odgovor studenta adekvatan.

Fokusiraj se na stavke koje se poklapaju sa studentovim odgovorom.

\end{tcolorbox}

\begin{tcolorbox}[promptbox, title={Neutralna (podrazumevana) konfiguracija}]

Ti si profesor visokog obrazovanja i treba da oceniš odgovor studenta na postavljeno pitanje.

Dobijaš odgovor studenta i pitanje na koje je student odgovarao, kao i tačan odgovor.

\end{tcolorbox}

\begin{tcolorbox}[promptbox, title={Stroga konfiguracija}]

Ti si strog profesor visokoakademske institucije i treba da ocenis odgovore studenata.

Dobijas pitanje, tacan odgovor kao i odgovor studenta, i treba da proveris da li je odgovor studenta adekvatan.

Svako nepoklapanje se kaznjava sa oduzimanjem petine ukupnih poena.

Fokusiraj se na stavke koje se ne poklapaju sa studentovim odgovorom.

\end{tcolorbox}

Neutralna konfiguracija je namerno najkraća i ne uvodi dodatna pravila;
blaga i stroga usmeravaju pažnju modela na poklapanja odnosno na nepoklapanja, a stroga zadaje i meru kazne.

\section*{Prompt za generisanje referentnog odgovora (GRA režim)}

U prvom koraku GRA režima model iz nastavnog materijala konstruiše referentni odgovor sledećom instrukcijom.
Korak ne zavisi od strogosti; ona se primenjuje tek pri ocenjivanju, istim promptovima kao u PRA režimu.

\begin{tcolorbox}[promptbox]

Ti si profesor visokog obrazovanja i treba da odgovoriš na pitanje na osnovu relevantnog sadržaja iz udžbenika.

Generiši koncizan odgovor koji direktno odgovara na postavljeno pitanje.

Odgovor treba da bude precizan i dužine od 2 do 5 rečenica.

\end{tcolorbox}

\section*{Promptovi za razmatrani treći režim}

Instrukcije za varijantu sa nastavnim materijalom neposredno u promptu (odeljak~\ref{sec:rezim-udzbenik})
razlikuju se od PRA verzija samo u opisu priloženog materijala.
Kao ilustracija navedena je neutralna verzija.

\begin{tcolorbox}[promptbox, title={Neutralna konfiguracija, varijanta sa nastavnim materijalom}]

Ti si profesor visokog obrazovanja i treba da oceniš odgovor studenta na postavljeno pitanje.

Dobijaš odgovor studenta i pitanje na koje je student odgovarao, kao i udzbenik iz kojeg je pitanje postavljeno.

\end{tcolorbox}

\section*{Reprezentativni primeri ocenjivanja}

Naredni primeri preuzeti su iz rezultata ocenjivanja modelom GPT-5 u neutralnoj konfiguraciji PRA režima, sa kursa Programiranje 2.
Izabrani su tako da ilustruju oba smera odstupanja od nastavničke ocene.

\section*{Primer A: sistem ocenjuje strože od nastavnika}

\begin{tcolorbox}[promptbox]

\textbf{Pitanje}

Koliko tačkica ispisuje poziv \texttt{f(3)} za narednu funkciju \texttt{f}?

\begin{verbatim}
void f(unsigned n) {
    if (n > 0) {
        f(n-1);
        putchar('.');
        f(n-1);
    }
}
\end{verbatim}

\medskip
\textbf{Referentni odgovor nastavnika}
\nopagebreak

Ispisuje 7 tackica. $T(n) = 2 \cdot T(n-1) + 1$ sa baznim slučajem $T(0) = 0$.

\medskip
\begin{samepage}
\textbf{Studentski odgovor}

\begin{verbatim}
Ispisuje 6 tačaka.

f(3), 3>0
  f(2), 2>0
    f(1), 1>0
      f(0)
      .
      f(0)
    .
    f(1), 1>0
      f(0)
      .
      f(0)
  .
  f(2), 2>0
    f(1), 1>0
      f(0)
      .
      f(0)
    .
    f(1), 1>0
      f(0)
      .
      f(0)
\end{verbatim}
\end{samepage}

\medskip
\textbf{Ocene}
\nopagebreak

Nastavnik: 5/10 \hfill Sistem: 3/10

\medskip
\textbf{Obrazloženje sistema}
\nopagebreak

Netačan je krajnji rezultat: \texttt{f(3)} ispisuje 7 tačkica, ne 6.
Ispravno je $T(n) = 2 \cdot T(n-1) + 1$ sa $T(0) = 0$, pa je $T(1) = 1$, $T(2) = 3$, $T(3) = 7$.
U vašem tragu izvršavanja propušten je barem jedan ispis (npr.\ kod jedne od \texttt{f(1)} grana), što je dovelo do pogrešnog prebrojavanja.
Delimični poeni za pokušaj raščlanjivanja rekurzije, ali računanje i zaključak nisu tačni.

\end{tcolorbox}

Primer ilustruje tipičan izvor neslaganja: student je postupak razložio korektno, ali je konačan broj pogrešan.
Nastavnik je pokušaj nagradio polovinom poena, dok je model veću težinu dao netačnom rezultatu.

\section*{Primer B: sistem ocenjuje blaže od nastavnika}

\begin{tcolorbox}[promptbox]

\textbf{Pitanje}

Da li je pristup n-tom elementu efikasniji za niz, ili za povezanu listu?
Da li je pristup n-tom elementu efikasniji za niz, ili za dinamički niz?
Molim Vas obrazložite odgovor.

\medskip
\textbf{Referentni odgovor nastavnika}
\nopagebreak

Za niz. Za niz je $O(1)$, za listu $O(n)$.
Isto je. U oba slučaja je $O(1)$.

\medskip
\begin{samepage}
\textbf{Studentski odgovor}
\nopagebreak

pristup n-tom elementu je laksi za niz, posto se elementi u nizu dodaju uz numeraciju,
dok se kod listi pokazuje sa liste na listu bez znanja reda, sto bi znacilo da bi nam trebala petlja sa brojacem n dok ne dodjemo do n tog elemta liste

\end{samepage}

\medskip
\textbf{Ocene}
\nopagebreak

Nastavnik: 5/10 \hfill Sistem: 6/10

\medskip
\textbf{Obrazloženje sistema}
\nopagebreak

Student je tačno objasnio prvi deo: pristup n-tom elementu je efikasniji za niz nego za povezanu listu,
jer se u listi mora prolaziti element po element (linearno), dok je u nizu direktan pristup.
Međutim, nije odgovorio na drugi deo pitanja (niz vs.\ dinamički niz).
Tačan odgovor za drugi deo je da je pristup n-tom elementu isti ($O(1)$) i za statički i za dinamički niz,
jer oba koriste kontinuirani blok memorije i omogućavaju indeksirani pristup.

\end{tcolorbox}

Pitanje se sastoji od dva dela, a student je odgovorio samo na prvi.
Model je propust identifikovao i naveo tačan odgovor na izostavljeni deo,
ali je odgovoru ipak dodelio nešto više od polovine poena, dok je nastavnik dodelio tačno polovinu.
U oba primera obrazloženje sistema je sadržajno tačno i proverljivo i kada numerička ocena odstupa od nastavničke.
